\section{Low-Level Command Construction and Dynamic Macros}
\subsubsection{Dynamic Control Sequences with \texttt{\textbackslash csname...\textbackslash endcsname}}

The construct

\begin{lstlisting}[language=LaTeX]
\csname ba@caption@#2\endcsname
\end{lstlisting}

is used to create or access control sequence names dynamically at runtime. Instead of writing a fixed macro name such as \verb|\ba@caption@A|, the \verb|\csname...\endcsname| mechanism allows the macro name to be constructed from expandable content.



\paragraph{Syntax meaning}

The syntax

\begin{lstlisting}[language=LaTeX]
\csname <name components>\endcsname
\end{lstlisting}

instructs TeX to:

\begin{itemize}
\item concatenate all expandable tokens between \verb|\csname| and \verb|\endcsname|,
\item interpret the resulting string as a control sequence name,
\item and execute or define that control sequence.
\end{itemize}



\paragraph{Example of dynamic macro construction}

Given:

\begin{lstlisting}[language=LaTeX]
#2 = A
\end{lstlisting}

the expression

\begin{lstlisting}[language=LaTeX]
\csname ba@caption@A\endcsname
\end{lstlisting}

is formed, which corresponds to the macro:

\begin{lstlisting}[language=LaTeX]
\ba@caption@A
\end{lstlisting}

Thus, the mechanism allows the programmatic creation of structured namespaces.



\paragraph{Usage in definitions}

A common use is dynamic macro assignment:

\begin{lstlisting}[language=LaTeX]
\expandafter\gdef\csname ba@caption@#2\endcsname{#1}
\end{lstlisting}

This defines a macro whose name depends on the argument \verb|#2|, enabling storage of multiple independent values under systematically generated keys.



\paragraph{Why \texttt{\textbackslash expandafter} is needed}

Because \verb|\csname...\endcsname| must be fully constructed before assignment, TeX often requires:

\begin{lstlisting}[language=LaTeX]
\expandafter
\end{lstlisting}

to ensure that the name is expanded before \verb|\gdef| or \verb|\def| is applied.



\paragraph{Conceptual interpretation}

The construct can be understood as a form of \emph{string-to-macro conversion}:

\begin{itemize}
\item Input: a string-like token sequence (e.g.\ \texttt{ba@caption@A})
\item Output: an actual TeX control sequence (\verb|\ba@caption@A|)
\end{itemize}

This mechanism is fundamental for implementing dictionaries, registries, and indexed data structures in classical \TeX{} programming, where no native data containers exist.



\paragraph{Role in the bifurcation atlas system}

In the context of the bifurcation atlas package, \verb|\csname...\endcsname| is used to:

\begin{itemize}
\item store region-specific captions,
\item create indexed color variables,
\item build hierarchical naming structures such as \texttt{ba@c1}, \texttt{ba@c2}, etc.
\end{itemize}

It is therefore a core building block for simulating structured data inside the TeX macro system.